\documentclass[11pt]{article}
\usepackage[a4paper,margin=25mm]{geometry}
\usepackage{amsmath,amssymb,amsthm,booktabs,tabularx,longtable,array}
\usepackage[hidelinks]{hyperref}
\newtheorem{theorem}{Theorem}
\newtheorem{corollary}[theorem]{Corollary}
\newcommand{\Q}{\mathbb{Q}}
\newcommand{\coeff}[1]{[x^{#1}]}
\title{A ratio correction and sibling identities for Mathar's Catalan-transform recurrence}
\author{Haobo Ma\\
  \small ChronoAI Pte Ltd, Singapore
  \and Wenlin Zhang\\
  \small National University of Singapore\\
  \small ChronoAI Pte Ltd, Singapore}
\date{17 November 2026}
\hypersetup{pdftitle={A ratio correction and sibling identities for Mathar's Catalan-transform recurrence},
  pdfauthor={Haobo Ma; Wenlin Zhang}}
\begin{document}
\maketitle
\begin{abstract}
Using Mathar's root-generating-function theorem (2021, Theorem~5
\cite{mathar}), we give an alternative coefficient proof of his conjecture
in OEIS A157100, revision~14, FORMULA line~11 (5~February 2015), for
every $n\geq4$.
We correct the generating-function ratio accompanying its partial-sum
description and extend the calculation to the adjacent sequence A025262:
its linear recurrence implies Somos's bilinear identity by elimination of
the index. The same radical identity proves Mathar's six-term conjecture
for A025268 on its fully defined domain $n\geq6$.
We also derive the related coefficient formulas, shifted generating
functions, and the leading asymptotic constants.
\end{abstract}

\section{The recurrence}

Let $C_k=\binom{2k}{k}/(k+1)$ and
$C(z)=\sum_{k\geq0}C_kz^k$, so $C(z)=1+zC(z)^2$.
All power series below belong to $\Q[[x]]$, and differentiation is formal.
Define the Catalan transform introduced by Paul Barry in A157100
\cite{oeisa} by
\begin{equation}\label{eq:A}
 A(x)=(1+x)C\bigl(x(1-x^2)\bigr)=\sum_{n\geq0}a_nx^n.
\end{equation}
Mathar's FORMULA assertion is explicitly labelled ``Conjecture'' at
\href{https://raw.githubusercontent.com/oeis/oeisdata/57bbd254e161585873f56950ff19ba6dc4aeb66c/seq/A157/A157100.seq\#L11}
{revision~14, line~11}. It states that these coefficients satisfy the displayed
five-term recurrence. We prove it on the full domain where all five
indices are nonnegative. The source gives offset zero but omits the
onset; $n\geq4$ is this interpretation, rather than an additional quoted
condition.

Algebraic generating functions can be converted into polynomial-coefficient
recurrences by differential equations; Ekhad, Yang and Zeilberger
\cite{eyz} describe this method for Mathar's OEIS conjectures.
Here the relevant general result is Mathar's Theorem~5
\cite[\S4.2, equations (41)--(45)]{mathar}, with
$p=1-4x+4x^3$, $q=2x(1-x)$, $w=1$, $v=-1$ and $r=2$.
The proof below makes this specialization explicit and determines the
initial values and the precise onset directly.

\begin{theorem}\label{thm:main}
The sequence defined by \eqref{eq:A} starts with
$(a_0,a_1,a_2,a_3)=(1,2,3,6)$ and satisfies, for every integer $n\geq4$,
\begin{equation}\label{eq:rec}
\begin{split}
 (n+1)a_n+(-5n+1)a_{n-1}
 &+2(2n-1)a_{n-2}+2(2n-7)a_{n-3}\\
 &+2(-2n+7)a_{n-4}=0.
\end{split}
\end{equation}
These four initial values and \eqref{eq:rec} determine the sequence uniquely.
\end{theorem}

\begin{proof}
Put $t=x(1-x^2)$, $q_0=x(1-x)$ and $D=1-4x+4x^3$.
Since $t(0)=0$, substitution into the Catalan equation is defined.
The series $s=1-2tC(t)=1-2q_0A$ has constant coefficient $1$ and satisfies
$s^2=1-4t=D$. Thus it is the unique square root of $D$ with
constant coefficient $1$, and
\begin{equation}\label{eq:radical}
 2q_0A=1-s,\qquad D s'=\tfrac12D's.
\end{equation}
For the second identity, differentiate $s^2=D$ and multiply by $s$.
Differentiating the first identity and multiplying by $D$ now gives
\[
 2q_0D A'+2q_0'D A=-\tfrac12D's
                 =-\tfrac12D'+D'q_0A.
\]
Consequently,
\begin{equation}\label{eq:ode}
 P(x)\,xA'(x)+Q(x)\,A(x)=1-3x^2,
\end{equation}
where direct polynomial multiplication yields
\begin{align*}
 P&=(1-x)D=1-5x+4x^2+4x^3-4x^4,\\
 Q&=Dq_0'-\tfrac12D'q_0
       =1-4x+6x^2-2x^3-2x^4.
\end{align*}
This calculation takes place in formal series; it requires no analytic
convergence or division by $x$.

For $n\geq4$ and $0\leq j\leq4$,
$\coeff{n}x^j(xA')=(n-j)a_{n-j}$ and
$\coeff{n}x^jA=a_{n-j}$. The coefficient of $a_{n-j}$ on the left of
\eqref{eq:ode} is therefore $(n-j)P_j+Q_j$, where
$P_j=\coeff{j}P$ and $Q_j=\coeff{j}Q$:
\[
\begin{array}{c|rrrrr}
 j&0&1&2&3&4\\ \hline
 P_j&1&-5&4&4&-4\\
 Q_j&1&-4&6&-2&-2\\
 (n-j)P_j+Q_j&n+1&-5n+1&4n-2&4n-14&-4n+14
\end{array}
\]
The right side has coefficient zero for these $n$, proving
\eqref{eq:rec}. Expansion of \eqref{eq:A} through degree three gives
the stated initial values. Finally $n+1\ne0$ for $n\geq4$, so each
successive term is uniquely determined.
\end{proof}

For clarity about small indices, extend $a_j=0$ for $j<0$.
The left side of \eqref{eq:rec} then equals
$\mathbf1_{n=0}-3\mathbf1_{n=2}$ for every $n\geq0$.
Thus the zero-extended homogeneous formula fails at $n=0$ and $n=2$;
this does not affect the five nonnegative-index assertion.

\section{Coefficient formulas and the partial-sum relation}

The adjacent formulas in the source follow without another differential
equation argument. Write $D=1-4x+4x^3$ and let $s=\sqrt D$
denote its square root with constant coefficient one. Set
\[
 B(x)=\frac{1-s}{2}=x(1-x)A(x)=\sum_{m\geq1}b_mx^m.
\]
This is the generating function of the offset-one sequence A025262
\cite{oeisb}. In particular $b_1=b_2=b_3=1$.

\begin{corollary}\label{cor:coeff}
For $n\geq0$ and $m\geq1$, with $\binom{k}{j}=0$ for $j>k$,
\begin{align}
 a_n&=\sum_{k=0}^n
 (-1)^{\binom{n-k}{2}}\binom{k}{\lfloor(n-k)/2\rfloor}C_k,
       \label{eq:asum}\\
 b_m&=\sum_{j=0}^{\lfloor(m-1)/2\rfloor}
 (-1)^j\binom{m-2j}{j}C_{m-2j-1},\label{eq:bsum}\\
 a_n&=\sum_{m=1}^{n+1}b_m.\label{eq:partial}
\end{align}
The coefficients $b_m$ are positive, and $a_n$ is positive and strictly
increasing. Moreover, for $m\geq4$,
\begin{equation}\label{eq:brec}
 m b_m=2(2m-3)b_{m-1}-2(2m-9)b_{m-3}.
\end{equation}
\end{corollary}

\begin{proof}
In \eqref{eq:A}, expand
$C(t)=\sum_{k\geq0}C_kx^k\sum_{j=0}^k(-1)^j\binom{k}{j}x^{2j}$.
A contribution to $a_n$ has $n-k=2j$ or $2j+1$, according to
which term of $1+x$ is used. In either case
$j=\lfloor(n-k)/2\rfloor$ and
$\binom{n-k}{2}\equiv j\pmod2$. This proves \eqref{eq:asum}.
Likewise $B=tC(t)=\sum_{k\geq1}C_{k-1}x^k(1-x^2)^k$;
putting $k=m-2j$ gives \eqref{eq:bsum}. Terms with $j>k$ vanish
by the binomial convention.

Since $B/x$ is a power series, the identity
$A=(B/x)/(1-x)$ proves \eqref{eq:partial}, including its upper
index $n+1$. Also $B-B^2=x-x^3$, whence
$b_m=\sum_{i=1}^{m-1}b_ib_{m-i}$ for $m\geq4$.
Starting with $b_1=b_2=b_3=1$, induction proves positivity;
\eqref{eq:partial} gives $a_0=1$ and
$a_n-a_{n-1}=b_{n+1}>0$ for $n\geq1$.
Finally \eqref{eq:radical} implies
\[
 D B'+(2-6x^2)B=1-3x^2.
\]
The coefficient of $x^{m-1}$ is zero for $m\geq4$ and is
$m b_m+(-4m+6)b_{m-1}+(4m-18)b_{m-3}$, proving
\eqref{eq:brec}.
\end{proof}

Equation \eqref{eq:bsum} is Kruchinin's 2013 formula in A025262;
\eqref{eq:brec} is Kotesovec's 2015 recurrence there.
The 27~June 2022 comment by Irvine and Mathar in A157100
correctly states the partial-sum relation, but the printed
generating-function ratio has its direction reversed. With the source's
offsets it is
\[
 \frac{A(x)}{B(x)}=\frac1{x(1-x)},\qquad
 \frac{B(x)}{A(x)}=x(1-x),
\]
where the first ratio is interpreted in $\Q((x))$.
For the zero-based series $\widetilde B=B/x$, the ratio is instead
$A/\widetilde B=1/(1-x)$.

There is also a short route from these earlier formulas to the main
recurrence. The partial-sum relation gives
$b_{n+1}=a_n-a_{n-1}$. Substitution into Kotesovec's recurrence
at $m=n+1$ gives
\[
 (n+1)(a_n-a_{n-1})-(4n-2)(a_{n-1}-a_{n-2})
 +(4n-14)(a_{n-3}-a_{n-4})=0,
\]
which expands to \eqref{eq:rec} for $n\geq4$.

\section{Further consequences of the same radical}

\subsection{Somos's bilinear identity}

The following is the formula signed by Somos on 18~January 2015 in
A025262, FORMULA line~19, and in A056010 \cite{oeisb,oeisl}.

\begin{corollary}\label{cor:bilinear}
For every $n\geq1$,
\begin{equation}\label{eq:bilinear}
\begin{split}
0={}&b_n(16b_{n+1}-64b_{n+3}+22b_{n+4})
       +b_{n+1}(32b_{n+2}-14b_{n+3})\\
   &+b_{n+2}(16b_{n+3}-10b_{n+4})
       +b_{n+3}(2b_{n+3}+b_{n+4}).
\end{split}
\end{equation}
\end{corollary}
\begin{proof}
Write $u_j=b_{n+j}$ for $0\leq j\leq4$. Applying
\eqref{eq:brec} at $n+3$ and $n+4$ gives
\[
 n(u_3-4u_2+4u_0)=-3u_3+6u_2+6u_0,
 \qquad
 n(u_4-4u_3+4u_1)=-4u_4+10u_3+2u_1.
\]
Multiply the first right side by $u_4-4u_3+4u_1$, subtract
the second right side times $u_3-4u_2+4u_0$, and expand.
The result is exactly \eqref{eq:bilinear}. No division is required,
so vanishing factors cause no exceptional cases.
\end{proof}

\subsection{Shifted coefficients and row sums}

\begin{corollary}\label{cor:shifts}
Let $c_n=b_{n+2}$ for $n\geq0$ and
$L(x)=\sum_{n\geq0}c_nx^n$. Then
\begin{equation}\label{eq:L}
 L=\frac{B-x}{x^2}=\frac{1-2x-s}{2x^2}
  =\frac{2(1-x)}{1-2x+s},\qquad
 x+L=(1+xL)^2.
\end{equation}
Equivalently,
\begin{equation}\label{eq:Lcat}
\begin{split}
 L&=\frac{1-x}{1-2x}
 C\left(\frac{x^2(1-x)}{(1-2x)^2}\right),\\
 c_n&=\sum_{k=0}^{\lfloor n/2\rfloor}C_k
 \sum_{i=0}^{\min(k+1,n-2k)}
 (-1)^i\binom{k+1}{i}\binom{n-i}{n-2k-i}2^{n-2k-i}.
\end{split}
\end{equation}
Consequently $c_0=c_1=1$ and
$c_n=2c_{n-1}+\sum_{i=0}^{n-2}c_i c_{n-2-i}$ for $n\geq2$.
These are the coefficients of A056010, and \eqref{eq:bilinear}
holds with $b_j$ replaced by $c_j$ for every $n\geq0$.
The row sums of A176703, indexed from row zero, are $b_{n+1}$.
\end{corollary}
\begin{proof}
Substitute $B=x+x^2L$ into $B-B^2=x-x^3$ and cancel $x^2$;
this gives \eqref{eq:L}, its recurrence and the initial values.
Writing the same quadratic as
$(1-2x)L=(1-x)+x^2L^2$ and substituting the Catalan equation
gives \eqref{eq:Lcat}; the rationalized form in \eqref{eq:L}
follows from $(1-2x)^2-D=4x^2(1-x)$. Its sum follows by expanding
$\sum_{k\geq0}C_kx^{2k}(1-x)^{k+1}(1-2x)^{-2k-1}$.
The bilinear statement follows from Corollary~\ref{cor:bilinear}
with $n$ replaced by $n+2$.
For the row sums, define the row polynomials from A176703
\cite{oeist} by $f_0(y)=1$, $f_1(y)=y$, $f_2(y)=1$ and
$f_n(y)=\sum_{i=0}^{n-1}f_i(y)f_{n-1-i}(y)$ for $n\geq3$.
At $y=1$ this is exactly the defining convolution for $b_{n+1}$,
with the same three initial values; induction proves the assertion.
\end{proof}

These deductions prove the algebraic and shifted-coefficient statements
in A056010, not its lattice-path interpretation or its determinant
claims. The polynomial change $Y=-B$ also gives
$Y^2+Y=x^3-x$, the curve equation in Somos's 2023 comment in A025262.
This change alone asserts no relation to an elliptic divisibility sequence.

\subsection{The neighboring six-term conjecture}

The coefficient argument has a useful general form, which is a
specialization of Mathar's Theorem~5. If
$p(x)=\sum_{j=0}^d p_jx^j\in\Q[x]$ has $p(0)=1$,
degree $d$, and $U=(1-\sqrt p)/2=\sum_{n\geq1}u_nx^n$, with the
square root of constant coefficient one, then
\begin{equation}\label{eq:general}
 pU'-\tfrac12p'U=-\tfrac14p',\qquad
 \sum_{j=0}^{d}p_j(n-\tfrac{3j}2)u_{n-j}=0\quad(n>d).
\end{equation}
The first identity follows by differentiating $(1-2U)^2=p$;
the second is its coefficient of $x^{n-1}$: the two contributions
are $p_j(n-j)u_{n-j}$ and $-\tfrac12jp_ju_{n-j}$.

\begin{corollary}\label{cor:neighbor}
Let $d_1=d_2=d_3=d_4=1$ and
$d_n=\sum_{i=1}^{n-1}d_i d_{n-i}$ for $n\geq5$, as in A025268
\cite{oeisd}. Its generating function is
\[
 U(x)=\sum_{n\geq1}d_nx^n
      =\frac{1-\sqrt{1-4x+4x^3+8x^4}}2.
\]
For $n\geq5$ it satisfies Kotesovec's recurrence
\begin{equation}\label{eq:shortd}
 R_n:=nd_n-(4n-6)d_{n-1}+(4n-18)d_{n-3}
                      +8(n-6)d_{n-4}=0.
\end{equation}
For $n\geq6$ it satisfies Mathar's conjecture, signed
25~January 2015 in A025268,
\href{https://raw.githubusercontent.com/oeis/oeisdata/57bbd254e161585873f56950ff19ba6dc4aeb66c/seq/A025/A025268.seq\#L8}
{revision~27, FORMULA line~8}:
\begin{equation}\label{eq:longd}
\begin{split}
 nd_n+(n+1)d_{n-1}+10(-2n+5)d_{n-2}
 &+2(2n-9)d_{n-3}\\
 {}+2(14n-79)d_{n-4}+40(n-7)d_{n-5}&=0.
\end{split}
\end{equation}
\end{corollary}
\begin{proof}
The initial values and convolution give $U-U^2=x-x^3-2x^4$.
The root with constant coefficient zero is the displayed $U$.
Apply \eqref{eq:general} with $p=1-4x+4x^3+8x^4$ to obtain
\eqref{eq:shortd}. Expanding $R_n+5R_{n-1}$ gives
\eqref{eq:longd}. Both shorter identities hold when $n\geq6$,
which is also the full domain of six positive indices.
\end{proof}

This deduction uses the shorter recurrence posted by Kotesovec on the
same day as Mathar's conjecture; the identities are not presented as
new recurrence formulas. Both sources omit an onset. With $d_0=0$,
the left side of \eqref{eq:longd} at $n=5$ is $-40$, rather than zero;
the zero-extended formula at that index was not asserted in the source.

Adamson's two constructions follow from the same convolution.
For A025262 put $e_n=b_{n+1}$; for A025268 put $e_n=d_{n+1}$.
The initial blocks are three or four ones. Splitting off the last
term of the convolution gives, after that block,
\[
 e_n=e_{n-1}+\sum_{j=0}^{n-2}e_j e_{n-1-j},
\]
because $e_0=1$. This is precisely the source's $\Phi$ construction
with its zero-based indexing.

\subsection{Leading asymptotic constants}

\begin{corollary}\label{cor:asymptotic}
Let $\rho$ be the unique zero of $D=1-4x+4x^3$ in $(1/4,1/3)$.
Then
\begin{equation}\label{eq:asymptotic}
\begin{aligned}
 b_n&\sim\frac{\sqrt{3-8\rho}}{4\sqrt\pi}\rho^{-n}n^{-3/2},\\
 a_n&\sim\frac{\sqrt{3-8\rho}}{4\sqrt\pi\rho(1-\rho)}
                         \rho^{-n}n^{-3/2}.
\end{aligned}
\end{equation}
\end{corollary}
\begin{proof}
The three roots of $D$ lie in $(-2,-1)$, $(1/4,1/3)$ and
$(3/4,1)$, by sign changes in these disjoint intervals; thus
$\rho$ is the unique root of smallest modulus, and is simple.
Factoring at $\rho$ gives
\[
 \sqrt{D(x)}=\sqrt{-\rho D'(\rho)}(1-x/\rho)^{1/2}
                  \bigl(1+O(1-x/\rho)\bigr).
\]
The branch with value one at zero continues to a slit neighborhood
of the closed disc of radius $\rho$; the other two roots have larger
modulus. The square-root transfer theorem
\cite[Theorems~VI.1 and VI.4]{fs}, with
$\Gamma(-1/2)=-2\sqrt\pi$, applied to $B=(1-\sqrt D)/2$
gives the first equivalent. Since $D(\rho)=0$,
$-\rho D'(\rho)=3-8\rho>0$.
Multiplication by $1/[x(1-x)]$, analytic at $\rho$, gives the second.
\end{proof}

The first equivalent is Kotesovec's 5~June 2022 formula in A025262:
$\rho^{-1}=4-4\rho^2$, and the trigonometric expression printed
there lies in $(1/4,1/3)$ and solves $D(\rho)=0$ by the triple-angle
identity. The second equivalent is a consequence for the partial sums,
not a conjecture posed in the source.

\clearpage
\section{Coverage and verification}

{\small\setlength{\tabcolsep}{4pt}
\begin{longtable}{@{}p{0.34\textwidth}p{0.62\textwidth}@{}}
\toprule
Source statement & Coverage\\
\midrule
A157100 five-term conjecture & Reproved; prior proof framework: Mathar
(2021), Theorem~5; explicit coefficients and onset in
Theorem~\ref{thm:main}.\\
A157100 binomial formula & Proved in \eqref{eq:asum}.\\
A157100 partial sums and ratio comment & Partial sums proved in
\eqref{eq:partial}; printed ratio refuted-and-corrected by
$A/B=1/[x(1-x)]$. The partial-sum identification is already
in the source's 2022 comment.\\
A025262 quadratic equation and convolution & Reproved in the proof of
Corollary~\ref{cor:coeff}; the quadratic formula is credited to Somos (2000).\\
A025262 binomial formula and linear recurrence & Reproved in
\eqref{eq:bsum} and \eqref{eq:brec}; prior formulas of Kruchinin (2013)
and Kotesovec (2015).\\
A025262 and A056010 bilinear formula & Proved in
Corollaries~\ref{cor:bilinear}--\ref{cor:shifts}; formula due to
Somos (2015), obtained from Kotesovec's linear recurrence.\\
A025262 and A025268 Adamson constructions & Proved: putting
$e_n=b_{n+1}$ or $d_{n+1}$ gives
$e_n=e_{n-1}+\sum_{j=0}^{n-2}e_j e_{n-1-j}$ for
$n\geq3$ or $n\geq4$, after three or four initial ones, respectively.\\
A176703 row sums; A056010 shifts and g.f. & Proved in
Corollary~\ref{cor:shifts}; the source's algebraic identities are
rederived with their offsets.\\
A025268 six-term conjecture & Proved for $n\geq6$ in
Corollary~\ref{cor:neighbor}, using the shorter Kotesovec (2015)
formula; its general mechanism is Mathar (2021).\\
A025262 asymptotic; A157100 consequence & Prior formula of Kotesovec
(2022) rederived in Corollary~\ref{cor:asymptotic}; the A157100
consequence is not posed in the source.\\
A025262 curve equation & Proved by $Y=-B$; no elliptic-divisibility
claim follows from that substitution alone.\\
A025262 ascent sequences & Prior proof cited: Callan (2019),
\S6, Theorem~5 \cite{callan}, including the offset $b_{n+1}$
for objects of length $n$.\\
A056010 path interpretation and generalized Somos-5 number wall &
Not covered by these coefficient proofs. Its alternative Catalan
form and binomial sum are proved in \eqref{eq:Lcat}.\\
A157100 Hankel transform; A157101 and its linked Somos variants &
Not covered by the coefficient argument. Barry (2019), Example~45
\cite{oeish,barry}, gives related A157101 Hankel context; no new assertion
about these determinants is made. Yura (2014), \S3,
and Barry (2023), \S8 \cite{yura,barry2023}, give the related
A025262/A056010 elliptic and Hankel constructions.\\
\bottomrule
\end{longtable}}

The pinned Lean source
\href{https://github.com/the-omega-institute/trureturing/blob/962b1e14510f845c23345d34b4fb8205b2ddc2bd/D5/S1/Recurrence/Algebraic/MatharA157100FiveTermRecurrence.lean}
{\texttt{MatharA157100FiveTermRecurrence.lean}}
defines \texttt{A} by \eqref{eq:A}, defines \texttt{a} as its rational
coefficients, and its declaration \texttt{result} proves
\texttt{claim}, exactly \eqref{eq:rec} for $n\geq4$.
The declaration \texttt{cat\_ode} in the pinned dependency
\href{https://github.com/the-omega-institute/trureturing/blob/962b1e14510f845c23345d34b4fb8205b2ddc2bd/D5/S1/Recurrence/A308815DFiniteRecurrence.lean}
{\texttt{A308815DFiniteRecurrence.lean}}
proves $x(1-4x)C'+(1-2x)C=1$; its \texttt{euler\_coeff}
proves $[x^n](xF')=n[x^n]F$.
The initial values and uniqueness in Theorem~\ref{thm:main},
Corollaries~\ref{cor:coeff}--\ref{cor:asymptotic}, the ratio correction
and the general identity \eqref{eq:general} have written proofs;
they are separate from the statement of \texttt{result}.
The accompanying \texttt{verify.py} uses integer arithmetic to compare
the Catalan substitution, quadratic convolution and recurrence,
and checks the displayed bilinear, shifted, row-sum and A025268 identities.
Its ranges and exact polynomial checks are specified in the accompanying
\texttt{README.md} and \texttt{verification.json}.
These are finite checks alongside the all-index proofs above.

\begin{thebibliography}{10}
\small
\bibitem{oeisa}
OEIS Foundation Inc., \emph{A157100}, entry by Paul Barry,
five-term conjecture by R.~J.~Mathar (5~February 2015), revision~14
(22~September 2025), FORMULA; ratio comment (27~June 2022).
\url{https://oeis.org/A157100}.
\bibitem{mathar}
R.~J.~Mathar, \emph{P-finite recurrences from generating functions with
roots of polynomials}, arXiv:2109.02112v1 (2021), \S4.2, Theorem~5.
\url{https://arxiv.org/abs/2109.02112v1}.
\bibitem{eyz}
S.~B.~Ekhad, M.~Yang and D.~Zeilberger,
\emph{Automated proofs of many conjectured recurrences in the OEIS
made by R.~J.~Mathar}, arXiv:1707.04654v2 (2017).
\url{https://arxiv.org/abs/1707.04654v2}.
\bibitem{oeisb}
OEIS Foundation Inc., \emph{A025262}, entry by Clark Kimberling,
revision~76 (5~November 2025), FORMULA; contributions of M.~Somos
(2000), V.~Kruchinin (2013), V.~Kotesovec (2015).
\url{https://oeis.org/A025262}.
\bibitem{oeish}
OEIS Foundation Inc., \emph{A157101}, entry by Paul Barry,
revision~13 (22~September 2025).
\url{https://oeis.org/A157101}.
\bibitem{oeisd}
OEIS Foundation Inc., \emph{A025268}, entry by Clark Kimberling,
revision~27 (4~February 2025), FORMULA; conjecture of R.~J.~Mathar
and generating function and recurrence of V.~Kotesovec (25~January 2015).
\url{https://oeis.org/A025268}.
\bibitem{oeist}
OEIS Foundation Inc., \emph{A176703}, entry by Roger L.~Bagula,
revision~13 (26~January 2024), FORMULA and row sums.
\url{https://oeis.org/A176703}.
\bibitem{oeisl}
OEIS Foundation Inc., \emph{A056010}, entry by Michael Somos,
revision~58 (23~July 2026), FORMULA.
\url{https://oeis.org/A056010}.
\bibitem{fs}
P.~Flajolet and R.~Sedgewick, \emph{Analytic Combinatorics},
Cambridge University Press, 2009, Chapter~VI,
Theorems~VI.1 and VI.4.
\url{https://algo.inria.fr/flajolet/Publications/book.pdf}.
\bibitem{barry}
P.~Barry, \emph{Generalized Catalan recurrences, Riordan arrays,
elliptic curves, and orthogonal polynomials}, arXiv:1910.00875v1
(2019), Example~45, p.~33.
\url{https://arxiv.org/abs/1910.00875v1}.
\bibitem{callan}
D.~Callan, \emph{On Ascent, Repetition and Descent Sequences},
arXiv:1911.02209v1 (2019), \S6, Theorem~5.
\url{https://arxiv.org/abs/1911.02209v1}.
\bibitem{yura}
F.~Yura, \emph{Hankel Determinant Solution for Elliptic Sequence},
arXiv:1411.6972v2 (2014), \S3.
\url{https://arxiv.org/abs/1411.6972v2}.
\bibitem{barry2023}
P.~Barry, \emph{Integer sequences from elliptic curves},
arXiv:2306.05025v2 (2023), \S8.
\url{https://arxiv.org/abs/2306.05025v2}.
\end{thebibliography}

\noindent\small\textcopyright\ 2026 Haobo Ma and Wenlin Zhang.
Licensed under Creative Commons Attribution 4.0 International
(CC BY 4.0).
\end{document}
